\section{商业模式创新：outcome-based pricing 与激励对齐}

讲者回顾了软件定价演化：从盒装授权、SaaS 订阅和按用量计费，再到按结果定价（outcome-based pricing）。Sierra 的做法是“问题解决才计费”，例如完成退换货、触发有效订单更新或完成其他可核验服务动作后再收费。

这种定价把平台方和企业客户绑定到同一目标函数：客户只有在节省成本或提升转化时才愿意持续投入，平台也必须把注意力放在“稳定解决问题”而非“最大化 token 消耗”。

\begin{knowledgebox}{outcome-based pricing 的技术前提}
要做 outcome-based，必须先解决 outcome verification：
\begin{itemize}[leftmargin=1.5em]
    \item 动作可验证：例如订单状态已更新、工单已关闭、退款已提交；
    \item 证据可追溯：可回放调用链、参数、策略版本；
    \item 归因可计算：成功来自 Agent 自动完成，而不是人工接管。
\end{itemize}
\end{knowledgebox}

\begin{importantbox}{实践启发}
如果团队还无法建立“可验证动作闭环”，就不宜直接承诺 outcome-based 商业模式。应该先完成事件模型、审计日志、状态变更对账，再谈计费创新。
\end{importantbox}

\subsection{本章小结}

outcome-based pricing 并非销售包装，而是架构承诺。它倒逼团队建立强可观测、强可归因、强可验证的 Agent 系统，这正是企业级落地的关键护城河。

